%Math Stuff
\usepackage{amsmath} % many display options for math modes and such (VERY IMPORTANT)
\usepackage{physics} % useful math symbols and macros (VERY IMPORTANT)
\usepackage{amssymb} % heck ton of math symbols and fonts (VERY IMPORTANT)
\usepackage{amsfonts} % Additional fonts, math symbols, and options for existing fonts (IMPORTANT)
\usepackage{mathtools} % lots of cool mathtools oh wait that is the name (IMPORTANT)
\usepackage{siunitx} % allows you to quickly define units within math mode (IMPORTANT)
\usepackage{array} % allows you to write piece-wise functions, matrices, and other cool things (IMPORTANT)
\usepackage{txfonts} % defines times new roman as default text font and provides supporting math symbols
\usepackage{braket} % allows the use of detailed Dirac braket notation
\usepackage{cancel} % used to cancel terms (with X or diagonal line) while writing math (useful for assignments)
\usepackage{ulem} % underlines, typewriter font (for comp sci somewhat), and striking out (required for cancel package)
%\usepackage{halloweenmath} % spooky
\usepackage{empheq} %Extension of amsmath for boxing equations or answers


%Formatting Stuff
\usepackage[utf8]{inputenc} % the typesetting rules
\usepackage[left = 22mm, right = 22mm]{geometry} % flexible and easy interface to change page dimensions
\usepackage{graphicx} % provides additional options for figures
\usepackage{float} % allows you to tell latex no really I want the figure HERE (with H)
\usepackage{color} % provides coloring capabilities to everything
\usepackage{fancyhdr} % making fancy headers (obviously)
\usepackage{changepage} % allows to change page formatting in the middle of a document (rather than the same throughout) 
\usepackage{enumitem} % better control over enumerate and itemize
\usepackage[labelfont=bf]{caption} % altering options for captions
\usepackage{sidecap} % allows sideways captions for figures and tables
\usepackage{soul,xcolor} % can enhance striking, underlining, and highlighting (in non-math mode)


%Stuff for Tables
\usepackage{booktabs} % allows weird lines within tables
\usepackage{multirow}
\usepackage{multicol} % allows for multiple columns and rows that collapse into single columns
\usepackage{longtable} % allows a table to run over multiple pages
\usepackage{colortbl} % colourful tables!

%Stuff for Citing
%\usepackage{apacite} % helps you apa cite
%\usepackage[numbers]{natbib} % implements both author-year and numbered references
\usepackage{url} % allows options for url links
\usepackage[colorlinks, citecolor = black, urlcolor = black, linkcolor = black]{hyperref} %references (citation) options and hyperlink integration


%Chemistry Stuff
\usepackage[version=4]{mhchem} % write chemical equations within math mode
\usepackage{chemfig} % draw simple and complex chemical structures in latex

%Elijah Adams Created Commands
\newcommand{\Q}{\mathbb{Q}} %Mathematical sets
\newcommand{\R}{\mathbb{R}}
\newcommand{\C}{\mathbb{C}}
\newcommand{\Z}{\mathbb{Z}} 
\newcommand{\Curl}{\vec{\nabla}\times} % Good looking curl and div and laplace (because physics package ones look like shit)
\newcommand{\Div}{\vec{\nabla}\cdot}
\newcommand{\laplace}{\vec{\nabla}^2}
\newcommand{\yas}{\frac{1}{4\pi\epsilon_0}} % Useful for E&M stuff
\newcommand{\mas}{\frac{\mu_0}{4\pi}}
\newcommand{\ppdv}[2]{\frac{\partial^2 #1}{\partial #2^2}} % Twice partial derivatives

\setlength{\parindent}{4mm}
\definecolor{purple}{rgb}{0.84,0.71,0.85} % defining a colour for later, each rgb value is a ratio of r/255, g/255 and b/255 because 255 is the maximum value of rgb numbers
\newcolumntype{d}{>{\columncolor{purple}}c} 

